\section{Introduction}
\label{sec:intro}
% Motivating gap: descriptor-based ML bottlenecked by O(10^3)-O(10^5) labeled datasets.
% Three-claim structure: Scale, Pipeline, Evaluation.
% Track-fit paragraph: explicit evaluative claims, scope, what the dataset cannot evaluate.
% Final paragraph: contributions list.

Machine learning and deep learning have upended the relationship between large datasets
and how we interpret, predict, and generate information.
Prior to this revolution, humans relied on low-dimensional embeddings and descriptors to 
parse and understand large swaths of data. Today, a wide spectrum of models operate on language~\citep{gpt3},
vision~\citep{yolo,sam}, and biological~\citep{alphafold2} data. In the chemical sciences, we have seen the decades long
development of descriptors and models heat up in the last 10 years. Cheminformatics via 
fingerprinting descriptors \citep{ecfp} and Morgan~\citep{morgan1965} coupling to tree-based supervised-learning algorithms\cite{ahneman2018}
have given way to advanced architectures and representations rooted in graph-based embeddings 
of chemical structures. These newer approaches include continuous-filter graph neural
networks~\citep{schnet}, directional message-passing architectures~\citep{dimenetpp},
and equivariant foundation models capable of stable atomistic simulation across
nearly the entire periodic table~\citep{macemp}. High-profile applications include
generative inorganic materials design with MatterGen~\citep{mattergen}, the OMat24
dataset and models for materials neural network potentials~\citep{omat24}, and 3D
pre-training on hundreds of millions of conformers for broad molecular property
prediction~\citep{unimol}.



Naturally, chemistry and material science, with their mature methods for generating simulated data via
Density Functional Theory (DFT) have led to the creation of many datasets. QM9, for example, is one of the
most popular datasets and covers a set of $\sim$134K organic molecules with up to 9 heavy atoms~\citep{qm9}.
Beyond this, SPICE~\citep{spice} and ANI-1x~\citep{ani1x} extended this coverage to biomolecular and reactive domains. OC20~\citep{oc20}, Materials Project~\citep{matproj}
and tmQM+~\citep{D5DD00220F} pushed the boundary towards inorganic materials and transition-metal complexes. Alongside datasets, benchmarking suites such as
MoleculeNet~\citep{moleculenet} and MOSES~\citep{moses} have standardized
evaluation across property prediction and molecular generation tasks. The recent OMol25~\citep{levine2026openmolecules2025omol25} is notable among molecular MLIP datasets due to 
its size, diversity, and high-level of accuracy. It pushes beyond 100M structures, 
80+ elements, and many chemical verticals to provide the community with a rich repository for 
training neural network potentials. To date, the team has shared the wavefunction, density matrices, 
and orca.out files of 4M of those datapoints. This additional information serves as the basis for this work 
where we post-processed OMol25 wavefunctions to yield a new dataset of descriptors. 
We name the result \textbf{OMol-Descriptors-4M} and contribute: (i) the dataset
itself, (ii) the open generation pipeline, and (iii) an evaluation protocol with
composition-ordered splits, five held-out stress-test suites, and quantified
per-descriptor noise floors.
